\documentclass[11pt]{article}
\usepackage[a4paper,margin=26mm]{geometry}
\usepackage{amsmath,amssymb}
\usepackage[hidelinks]{hyperref}
\setlength{\emergencystretch}{3em}
\title{Conjecture 00000000742: minimal components and minimal maps on $\mathbf Z_2$}
\author{lizaixi01, with OpenAI Codex assistance}
\date{22 October 2026}
\begin{document}
\maketitle
\section{Two readings of the statement}
The conjecture gives $p^{p-1}$ as the number of ``minimal decompositions
(minimal dynamical systems)'' of $1$-Lipschitz maps on $\mathbf Z_p$,
and ends with ``a count of minimal maps''. These phrases can refer to
different objects. At $p=2$ the proposed number is $2$. We refute both
the count of minimal components of a fixed map and the count of distinct
minimal maps. We do not identify different maps modulo topological
conjugacy: the source specifies no such equivalence relation.

A minimal closed subsystem means a nonempty closed forward-invariant
set $K$ such that every forward orbit from a point of $K$ is dense in
$K$. In the terminology of Fan--Liao \cite{FL}, their minimal components
are certain clopen minimal subsystems; an everywhere minimal system has
one component. Their general decomposition theorem concerns polynomials
of degree at least two. We use their terminology only, and prove the
affine examples below directly without invoking that theorem.

\section{A single map has one minimal component}
Let $T:\mathbf Z_2\to\mathbf Z_2$ be $T(x)=x+1$. Translation preserves
the $2$-adic distance, so $T$ is $1$-Lipschitz. For every $n\geq0$,
\[
 T^n(x)=x+n.
\]
The nonnegative integers are dense in $\mathbf Z_2$: every residue class
modulo $2^k$ has a nonnegative integer representative. Translation is a
homeomorphism. It follows that every forward orbit $\{x+n:n\geq0\}$ is
dense in the actual $2$-adic topology.

If $K$ is any nonempty closed minimal subsystem, choose $x\in K$.
Its orbit closure is both $K$ and $\mathbf Z_2$, so $K=\mathbf Z_2$.
Conversely, the whole space is a minimal subsystem, and is clopen.
Thus the collection of nonempty closed minimal subsystems is exactly
$\{\mathbf Z_2\}$. Its cardinality is $1$, whereas $2^{2-1}=2$.
This also proves the assertion for the smaller class of clopen minimal
components. No finite quotient is substituted for $\mathbf Z_2$.

\section{There are at least three distinct minimal maps}
For a unit $a\in\mathbf Z_2^\times$, let $T_a(x)=x+a$. This map is
$1$-Lipschitz and $T_a^n(x)=x+na$. Multiplication by a unit is a
continuous surjection of $\mathbf Z_2$ (in fact a homeomorphism).
Consequently it carries the dense nonnegative integers to a dense set,
and translation by $x$ shows that every forward $T_a$-orbit is dense.

The values $a=1,3,5$ are all units. Thus $T_1,T_3,T_5$ are minimal
$1$-Lipschitz maps. Their values at zero are distinct in $\mathbf Z_2$
(which has characteristic zero), so the maps themselves are distinct.
There are at least three such maps; there cannot be a bijection between
the set of these maps and a two-element set. These three examples are
mutually conjugate, which is irrelevant to counting actual maps. We make
no claim here about a count of conjugacy classes.

\section{Formalization and verification}
The Lean type \texttt{Z2} is Mathlib's actual \texttt{PadicInt 2}, with
its $2$-adic metric and topology. \texttt{odometer\_\allowbreak lipschitz} proves the
metric inequality, \texttt{odometer\_\allowbreak iterate} proves the orbit formula,
and \texttt{odometer\_\allowbreak dense\_\allowbreak orbit} uses the proved Mathlib theorem
\texttt{PadicInt.\allowbreak denseRange\_\allowbreak natCast}. The definition
\texttt{IsMinimal\allowbreak Component} states the closed-subsystem conditions above.
\texttt{odometer\_\allowbreak minimal\_\allowbreak iff} characterizes every such set as the whole
space; \texttt{minimal\_\allowbreak component\_\allowbreak count} proves its actual cardinality.

\texttt{IsMinimal\allowbreak Map} includes both the $1$-Lipschitz property and the
density of every forward orbit. \texttt{translation\_\allowbreak dense\_\allowbreak orbit}
proves density for every unit translation. \texttt{threeMinimal\allowbreak Maps}
constructs the three actual maps, and its injectivity is proved by
evaluation at zero. \texttt{not\_\allowbreak exactly\_\allowbreak two\_\allowbreak minimal\_\allowbreak maps} rules out
a bijection with the proposed finite set, without interpreting an
infinite cardinality as a natural number.

All eleven principal declarations were audited. The package pins Lean
4.33.0, Mathlib and all transitive dependencies. A separate project
build and official fresh-kernel replay accompany the source; no admitted
proof, native decision shortcut or extra axiom is used. The bilingual
source is preserved from upstream commit\par\noindent
\texttt{b355e42906e6eaf52bc7e9d005e0ef02d8aa7400}.\par
OpenAI Codex produced the proof and report. A separate supervising Codex
agent reviewed both readings and the full formal chain. This is AI
review, not human certification. The examples are classical and are
not claimed as new mathematical results.

\begin{thebibliography}{1}
\bibitem{FL} A.-H. Fan and L. Liao,
\emph{On minimal decomposition of $p$-adic polynomial dynamical systems},
Advances in Mathematics 228 (2011), 2116--2144, Introduction.
\url{https://arxiv.org/abs/1010.5583}.
\end{thebibliography}
\end{document}
